% ====================================================================
% §1 기호와 규약, 실험조건 매핑, 분극 (N0·N1) — 통합 절
% ====================================================================
\section{기호와 규약, 실험조건 매핑, 분극}\label{sec:notation}
계산 진행(그림~\ref{fig:spine})의 첫 두 노드는 전이 루프에 들어가기 전에 한 번만 실행된다 --- 실험조건을 모델
입력으로 환산하는 $\mathrm N0$ 와, 측정 전위에서 분극을 벗겨 내부 전위를 얻는 $\mathrm N1$ 이다. 이 절은 그 둘을
한데 묶어, 이후 모든 전이 계산이 딛고 설 기호$\cdot$부호 규약과 두 전위를 세운다.

\subsection{실험조건 매핑과 방향 부호}\label{sec:n0map}
계산의 출발점은 실험조건 $(V_\app,\text{direction},\text{c\_rate},Q_\cell,T)$ 을 받아 그 첫 일로 이를 모델
입력으로 환산하는 것이다 --- 이것이 노드 $\mathrm N0$ 다. 방향 문자열은 방향 부호 $\sigma_d$ 로,
C-rate 는 전류 크기 $|I|$ 로 바뀐다:
\begin{equation}
\sigma_d=\begin{cases}+1 & \text{방전(discharge)}\\ -1 & \text{충전(charge)}\end{cases},
\qquad
|I|=\text{c\_rate}\cdot Q_\cell\quad(\text{또는 }I_\mathrm{abs}\text{ 직접 지정}).
\label{eq:n0map}
\end{equation}
c\_rate 는 시간 역수이므로 $Q_\cell$ 과 시간 단위를 맞춰 대입한다([1/h]$\cdot$[A\,h]$\to$[A]; 정규화
$Q_\cell{=}1$ 에서는 무차원 배율). 방향 부호가 실제로 작용하는 자리는 셋뿐이다 --- 분극 부호($\mathrm N1$)$\cdot$분기 중심 선택($\mathrm N3$)$\cdot$꼬리($\mathrm N7\cdot\mathrm N8$:
방향별 전달계수 $\chi_d$ 가 꼬리 \emph{길이}를, 적분 방향(식~\eqref{eq:reversal})이 인과 기억의 진행 \emph{방향}을 가른다).
그 밖에서 평형 종 자체는 방향 불변이며(아래 절들에서 식으로 회수), 본 장이 뽑는 양은 $\dd Q/\dd V$ 한 곡선이다.

\textbf{전이 인덱스.} $j=1,\dots,N_p$($N_p\approx3$--$4$). 전이 $j$ 의 진행률 $\xi_j$(방전 시 $0\to1$ 로 증가, 곧 탈리튬화 진행),
용량 가중 $Q_j$. 박스 결과식은 $j$ 첨자를 단다. 점유율과 진행률은 $\theta_j=1-\xi_j$ 로 묶이며(점유 $\theta$ 가 $1\to0$
줄어드는 것이 진행 $\xi$ 가 $0\to1$ 느는 것), 유도에서 둘은 부호 한 번 차이로 오간다.

\textbf{부호 규약.} 전위는 흑연 음극 반쪽전지 전위(vs Li/Li$^+$)다. 방향 부호 $\sigma_d=\pm1$
(방전 $+1$/충전 $-1$)이며, 평형 진행률 $\xi_\eq=\mathrm{logistic}[\sigma_d(V-U)/w]$ 에서 방전
($\sigma_d=+1$)은 $V\!\uparrow\Rightarrow\xi_\eq\!\uparrow$(전위를 올리면 탈리튬화). 분기 중심 $+\tfrac12\sigma_d(\cdots)$
는 $\sigma_d$ 를 따라 방전$\cdot$충전이 반대로 갈리고, 꼬리는 충전에서 적분 방향이 반전된다. 이 규약은 \S\ref{sec:notation}--\S\ref{sec:sum} 전건에서 유지되고 \S\ref{sec:signcheck} 가 전수 검산하며,
Chapter 2(LCO)는 같은 규약을 전극-중립형으로 잇는다 --- $\sigma_d$ 슬롯의 물리는 ``탈리튬화 $=+1$''
그대로이고, 양극에서는 셀 라벨(충$\cdot$방전)이 식~\eqref{eq:lco-sigmaslot} 로 환산되어 들어온다(\S\ref{sec:lco-intro}). 유도 편의를 위해 \S\ref{sec:center}$\cdot$\S\ref{sec:hys}
의 $U_j$ 정의 관례에는 \emph{항상 $+1$ 인 고정 부호} $s$ 를 따로 두는데, 이는 방향에 따라 $\pm1$ 로 바뀌는 $\sigma_d$ 와
별개다 --- $s$ 는 Part 0 의 박스에서 명시적으로 유지되다가 본론의 결과-사슬 boxed 식에서 $s{=}1$ 로 대입돼 소멸한다.

\subsection{기호 마스터표와 흑연 staging 개관}\label{sec:notation-table}
표~\ref{tab:notation} 이 본론 전건에서 쓰는 기호$\cdot$단위$\cdot$의미의 마스터 대응이다. 실험조건 매핑 $\mathrm N0$
의 전제이자 이후 절의 참조 기준이며, 물리 기호와 구현 식별자의 대응은 부록~\ref{sec:appendix-code}(표~\ref{tab:symcode})에 따로 둔다.
이 표는 \emph{참조용} 조회표다 --- 각 기호는 해당 절에서 정식으로 도입$\cdot$유도되므로, 처음 읽을 때
전량을 새길 필요 없이 훑어 두고 본문 진행 중 기호를 만날 때 되짚는 용도로 쓰면 된다.

\begin{longtable}{p{0.15\textwidth} p{0.16\textwidth} p{0.605\textwidth}}
\caption{기호$\cdot$단위$\cdot$의미 대응표 --- N0 실험조건 매핑 전제.}\label{tab:notation}\\
\toprule 기호 & 단위 & 의미 \\ \midrule \endfirsthead
\toprule 기호 & 단위 & 의미 \\ \midrule \endhead
\multicolumn{3}{l}{\emph{— 실험조건$\cdot$방향(N0$\cdot$N1) —}}\\
$\sigma_d$ & --- & 방향 부호 — 방전 $+1$/충전 $-1$ \\
$s$ & --- & 유도 전용 고정 부호 — 항상 $+1$(\S\ref{sec:center}$\cdot$\S\ref{sec:hys} 의 $U_j$ 정의 관례); 방향에 따라 $\pm1$ 로 바뀌는 $\sigma_d$ 와 별개, 본론 결과-사슬 boxed 식에서 $s{=}1$ 대입돼 소멸(Part 0 의 박스는 $s$ 를 명시 유지) \\
$|I|$ & A & 전류 크기 $=$ c-rate$\,\cdot Q_\cell$ \\
$Q_\cell$ & C(또는 A\,h) & 셀 용량(전하); $q=Q/Q_\cell$ 환산$\cdot$$|I|$ 환산 공용 — $|I|$ 환산 시 c-rate 와 시간 단위 정합 \\
$T$ & K & 온도 — 스칼라(등온) 또는 $V_\app$ 길이 배열($T(V)$ 비등온) \\
$T_\rep$ & K & 비등온 시 전이당 상수 평가에 공통으로 쓰는 \emph{전 구간 평균} $\overline T$(전이 창별 평균이 아니라 진행 구간 전체의 단일 평균) — N3$\cdot$N7 대표조건 평가(분기 중심$\cdot$지연 길이) \\
$R_n$ & $\Omega$ & 음극 lumped 분극 저항 \\
$V_\app$ & V & 인가(측정) 전위 배열 \\
$V_n$ & V & 내부(열역학) 전위 $=V_\app-\sigma_d|I|R_n$ \\
\multicolumn{3}{l}{\emph{— 평형 중심$\cdot$폭$\cdot$진행률(N2$\cdot$N4$\cdot$N5) —}}\\
$U_j(T)$ & V & 전이 $j$ 평형 중심 전위 $=(-\Delta H_\rxn+T\Delta S_\rxn)/F$ \\
$\Delta H_{\rxn,j},\Delta S_{\rxn,j}$ & {\footnotesize J/mol;\,J/(mol\,K)} & 삽입 \emph{반응} 엔탈피$\cdot$엔트로피(평형 — \emph{활성화}와 다름) \\
$n_j$ & --- & 폭 다중도($w=n_jRT/F$; 폭 $w$ 를 직접 주면 $n=wF/RT$ 역산) \\
$w_j$ & V & 전이 폭 척도 $=n_jRT/F$ \\
$\xi_{\eq,j}$ & --- & 평형 진행률(logistic) \\
$Q_j$ & C & 전이 $j$ 용량 가중 \\
$C_\bg$ & C/V & 배경 미분용량(상수 또는 $V$ 함수) \\
\multicolumn{3}{l}{\emph{— 히스테리시스 분기(N3) —}}\\
$\Omega_j$ & J/mol & 정규용액(regular solution) 상호작용 에너지($>2RT$ 면 상분리$\cdot$히스) \\
$T_{c,j}$ & K & 히스 문턱 임계온도 $=\Omega_j/2R$(식~\eqref{eq:sm-thresh}) \\
$\gamma_j$ & --- & 분기 축소 인자 $0\le\gamma_j\le1$($0$ 이면 분기 없음) \\
$\Delta U_j^\hys$ & V & spinodal 상한 분기 gap \\
$h_{\eta,j}$ & --- & 부분 cycle 인자(기본 $1$) \\
$U_j^{\,d}$ & V & 방향별 분기 중심 \\
\multicolumn{3}{l}{\emph{— 동역학 꼬리(N7$\cdot$N8) —}}\\
$\chi$ & --- & 전이상태 분율 위치(전달 계수); 방향별 $\chi_d$ \\
$\chi_d$ & --- & 방향별 전달 계수 — 방전 $\chi$/충전 $1-\chi$ \\
$\mathcal A_j$, $\mathcal A$ & J/mol & affinity(구동력) — 첨자형 $\mathcal A_j=sF(V-U_j)$ 는 유도용
일반 구동력(\S\ref{sec:width}; N7 일반형은 $-\Omega_j(1-2\xi_j)$ 가 더해짐), 무첨자 $\mathcal A$ 는
꼬리 컷점에 동결한 값 $=\min(z_\mathrm{cut}n_jRT,\,A_\mathrm{cap}RT)$(\S\ref{sec:lag}) \\
$\Delta H_{a,j},\Delta S_{a,j}$ & {\footnotesize J/mol;\,J/(mol\,K)} & \emph{활성화} 엔탈피$\cdot$엔트로피(속도) \\
$\Delta H_{a,j}^\eff$ & J/mol & 유효 장벽 $=\Delta H_a-\chi_d\Omega$ \\
$L_{q,j}$ & --- & 용량축 지연 길이 \\
$L_{V,j}$ & V & 전압축 지연 길이 $=|\dd V/\dd q|_{q_a}L_{q,j}$ \\
$|\dd V/\dd q|_{q_a}$ & V & 컷점 OCV(open circuit voltage) 기울기($L_q\to L_V$ 환산) \\
$z_\mathrm{cut}$ & --- & 컷 affinity 의 $z=\mathcal A/(n_jRT)$(기본 $4.357$) \\
$A_\mathrm{cap}$ & --- & 컷 affinity 상한 배수(기본 $4.0$) \\
$\xi_{\mathrm{lag},j}$ & --- & 인과 기억 적분된 진행률 \\
\multicolumn{3}{l}{\emph{— 상수 —}}\\
$R,F,k_B,h$ & {\footnotesize J/(mol\,K), C/mol, J/K, J\,s} & 기체상수, Faraday, Boltzmann, Planck \\
\bottomrule
\end{longtable}

흑연 staging 의 갤러리 채움을 그림~\ref{fig:staging} 에 둔다 --- 각 stage 전이가 한 peak 를 남기고,
흑연 staging 네 전이의 초기값(문헌 기반\cite{dahn1991,ohzuku1993} --- 값의 확정과 tier 는
\S\ref{sec:sum} 표~\ref{tab:staging} 소관)이 이 네 전이에 대응한다. 이 그림이 서론에서
말한 ``상전이 하나 $=$ peak 하나''의 실제 모양이며, 아래 $\dd Q/\dd V$ 축은 평형 종 식~\eqref{eq:eqpeak}(\S\ref{sec:eqpeak})을
표~\ref{tab:staging} 의 $U_j,Q_j$ 와 \S\ref{sec:sum} 의 \emph{폭 폴백 값} $w_j$($0.020/0.016/0.014/0.012$ V ---
$n$ 입력 제거 시 활성화되는 전이별 대안 폭)로 평가한 것이다. 실제 초기 상태는 $n_j{=}1$ 이 우선이라 네 전이
균일 $w=RT/F$ 로 평가되며(\S\ref{sec:sum}, 그림~\ref{fig:sumcurve}), 여기서는 전이별 폭$\cdot$높이 차이를
드러내는 폴백 폭 쪽을 보였다.

\begin{figure}[t]
\centering
\begin{tikzpicture}[x=1.05cm,y=0.5cm,
  lab/.style={font=\scriptsize,align=center},
  sub/.style={font=\tiny,align=center,black!55},
  fill4/.style={fill=black!72},fillL/.style={fill=black!24}]
% ---- (top) five stage columns: graphene layers + Li-filled galleries ----
\def\colw{0.85}
\foreach \x/\name/\ratio in {0/{stage 4}/{1:4}, 1.4/{stage 3}/{1:3}, 2.8/{stage 2L}/{1:2L},
                              4.2/{stage 2}/{1:2}, 5.6/{stage 1}/{1:1}}{
  \draw[semithick] (\x,0)--(\x+\colw,0);
  \foreach \k in {1,2,3,4,5,6}{\draw[semithick] (\x,0.6*\k)--(\x+\colw,0.6*\k);}
  \node[lab,below] at (\x+\colw/2,-0.15) {\name};
  \node[sub,below] at (\x+\colw/2,-0.95) {(\ratio)};
}
% galleries (bands between layers) — explicit per column, periodicity = staging index
\foreach \k in {1,5}{\fill[black!72] (0,0.6*\k-0.42) rectangle (0+\colw,0.6*\k-0.18);}
\foreach \k in {1,4}{\fill[black!72] (1.4,0.6*\k-0.42) rectangle (1.4+\colw,0.6*\k-0.18);}
\foreach \k in {1,3,5}{\fill[black!24] (2.8,0.6*\k-0.42) rectangle (2.8+\colw,0.6*\k-0.18);}
\foreach \k in {1,3,5}{\fill[black!72] (4.2,0.6*\k-0.42) rectangle (4.2+\colw,0.6*\k-0.18);}
\foreach \k in {1,2,3,4,5,6}{\fill[black!72] (5.6,0.6*\k-0.42) rectangle (5.6+\colw,0.6*\k-0.18);}
\draw[-{Latex[length=1.8mm]},thick] (0,4.05)--(6.45,4.05) node[pos=0.5,above,font=\scriptsize] {lithiation (charge): $4\to3\to2\mathrm L\to2\to1$};
\draw[{Latex[length=1.8mm]}-,thick] (0,-1.85)--(6.45,-1.85) node[pos=0.5,below,font=\scriptsize] {delithiation (discharge): main direction, $\theta_j:1\to0$};
% ---- (bottom) real dQ/dV axis: eq:eqpeak evaluated at tab:staging initial values ----
% x maps V in [0.05,0.25] onto the same gallery width [0,6.45]: x=6.45*(0.25-V)/0.20 (V decreasing rightward)
% yshift chosen (real cm, independent of the y=0.5cm/unit gallery above) so the tallest
% peak label (2->1, local y~10.9) clears the gallery's bottom arrow+label with margin.
\begin{scope}[yshift=-4.15cm,y=0.21cm]
\draw[-{Latex[length=1.5mm]}] (-0.9,0) -- (7.6,0) node[right,font=\scriptsize] {$V$ [V]};
\draw[-{Latex[length=1.5mm]}] (-0.9,0) -- (-0.9,11.3) node[above,font=\scriptsize,align=center] {$\dd Q/\dd V$\\{\tiny[$Q_\cell/$V]}};
\foreach \xx/\vv in {0/0.25,1.6125/0.20,3.225/0.15,4.8375/0.10,6.45/0.05}{
  \draw (\xx,0.35)--(\xx,-0.35) node[below,font=\scriptsize] {\vv};}
\foreach \yy in {0,4,8}{\draw (-0.95,\yy)--(-0.85,\yy) node[left,font=\tiny] {\yy};}
% 4->3 (U=0.210, w=0.020, Q=0.10)
\draw[thick] plot[smooth] coordinates {(4.515,0.033) (4.3215,0.045) (4.128,0.06) (3.9345,0.08) (3.741,0.107) (3.5475,0.142) (3.354,0.188) (3.1605,0.247) (2.967,0.322) (2.7735,0.414) (2.58,0.525) (2.3865,0.653) (2.193,0.793) (1.9995,0.937) (1.806,1.069) (1.6125,1.175) (1.419,1.238) (1.2255,1.247) (1.032,1.201) (0.8385,1.109) (0.645,0.983) (0.4515,0.842) (0.258,0.699) (0.0645,0.566) (-0.129,0.449) (-0.3225,0.35) (-0.516,0.27)};
% 3->2L (U=0.140, w=0.016, Q=0.12)
\draw[densely dashed,thick] plot[smooth] coordinates {(6.063,0.056) (5.8695,0.082) (5.676,0.117) (5.4825,0.168) (5.289,0.24) (5.0955,0.339) (4.902,0.472) (4.7085,0.647) (4.515,0.865) (4.3215,1.119) (4.128,1.388) (3.9345,1.634) (3.741,1.811) (3.5475,1.875) (3.354,1.811) (3.1605,1.634) (2.967,1.388) (2.7735,1.119) (2.58,0.865) (2.3865,0.647) (2.193,0.472) (1.9995,0.339) (1.806,0.24) (1.6125,0.168) (1.419,0.117) (1.2255,0.082) (1.032,0.056)};
% 2L->2 (U=0.120, w=0.014, Q=0.25)
\draw[black!55,thick] plot[smooth] coordinates {(6.45,0.119) (6.2565,0.181) (6.063,0.275) (5.8695,0.415) (5.676,0.621) (5.4825,0.917) (5.289,1.33) (5.0955,1.875) (4.902,2.543) (4.7085,3.274) (4.515,3.94) (4.3215,4.374) (4.128,4.442) (3.9345,4.119) (3.741,3.511) (3.5475,2.785) (3.354,2.086) (3.1605,1.496) (2.967,1.041) (2.7735,0.708) (2.58,0.475) (2.3865,0.315) (2.193,0.208) (1.9995,0.137)};
% 2->1 (U=0.085, w=0.012, Q=0.50)
\draw[very thick] plot[smooth] coordinates {(7.224,0.301) (7.0305,0.491) (6.837,0.797) (6.6435,1.282) (6.45,2.029) (6.2565,3.133) (6.063,4.658) (5.8695,6.545) (5.676,8.502) (5.4825,9.977) (5.289,10.399) (5.0955,9.579) (4.902,7.872) (4.7085,5.888) (4.515,4.103) (4.3215,2.72) (4.128,1.745) (3.9345,1.096) (3.741,0.679) (3.5475,0.417)};
% legend, placed in a region genuinely clear of all four curves (x<2.6, y>6.8; verified
% against the coordinate lists above -- the tallest curve there, 4->3, peaks at 1.247)
\draw[thick] (0.15,10.85)--(0.75,10.85); \node[right,font=\tiny] at (0.8,10.85) {$4\to3$, $U{=}0.210$ V};
\draw[densely dashed,thick] (0.15,9.55)--(0.75,9.55); \node[right,font=\tiny] at (0.8,9.55) {$3\to2\mathrm L$, $U{=}0.140$ V};
\draw[black!55,thick] (0.15,8.25)--(0.75,8.25); \node[right,font=\tiny,black!55] at (0.8,8.25) {$2\mathrm L\to2$, $U{=}0.120$ V};
\draw[very thick] (0.15,6.95)--(0.75,6.95); \node[right,font=\tiny] at (0.8,6.95) {$2\to1$, $U{=}0.085$ V};
\end{scope}
\end{tikzpicture}
\caption{흑연 staging 의 갤러리 채움 도식(위)과 대응 $\dd Q/\dd V$ peak(아래, 식~\eqref{eq:eqpeak} 의
수치 평가 --- $\dd Q/\dd V=Q_j\xi_\eq(1-\xi_\eq)/w_j$, 표~\ref{tab:staging} 의 $U_j,Q_j$ 와 \S\ref{sec:sum} 폭 폴백 $w_j$).
위: 수평선이 그래핀 층, 음영 띠가 리튬이 든 층간 갤러리다. 채울수록 stage 번호가 줄어 stage 1($\mathrm{LiC_6}$,
완전 채움)에 닿고, 열 아래 비율(1:4$\to$1:1)이 각 stage 의 갤러리 점유 주기다. stage $2\mathrm L$ 은 면내 질서가
풀린 단계라 옅게 표시했다. 아래: 폭 폴백 값으로 평가한 네 전이 peak --- 높이 $Q_j/(4w_j)$ 는 전이마다 크게 다르다
($2\to1$ 이 가장 큼, $10.4\,Q_\cell/$V; $n_j{=}1$ 우선의 실제 초기 폭 $w{=}RT/F$ 에서는 $4.87$ --- \S\ref{sec:sum}). \emph{주의} --- 위 갤러리 열 간격은 범주형(등간격 스케치)이고 아래 $V$
축은 실제 전위 스케일(비등간격 --- $3\to2\mathrm L$ 과 $2\mathrm L\to2$ 은 겨우 $20$ mV 차이라 갤러리 열 간격에
선형 정렬되지 않는다)이라, 둘은 전이명(예 $2\to1$)으로만 대응시킨다. 본 장 방향은 방전(탈리튬화)이다.}
\label{fig:staging}
\end{figure}

이 골격의 두 번째 전극(LCO 양극) 일반화는 Chapter 2(\S\ref{sec:lco-intro})가 연다. 그에 앞서, $\mathrm N0$ 의
남은 한 걸음인 분극 환산($\mathrm N1$)을 다음 소절이 마친다.

\subsection{분극 --- 측정 전위에서 내부 전위로}\label{sec:pol}
계산의 첫 물리식은 분극을 벗기는 환산이다. 관측은 유한 전류에서 이루어지므로 측정 전위 $V_\app$ 에는 셀의
직렬 저항을 통과한 IR 분극이 실려 있고, 평형 열역학이 성립하는 전위는 그 분극을 벗긴 내부 전위 $V_n$ 이다.

\textbf{(a) 출발 --- 측정 전위와 내부 전위의 차.} 전류 크기 $|I|$ 가 lumped 저항 $R_n$ 을 통과하며 만드는 전위 강하는
크기 $|I|R_n$ 이고, 그 부호는 전류 방향이 정한다. 측정 전위는 내부 전위에 분극이 더해진 값이므로 한 방향(방전)에서
\begin{equation}
V_\app\;=\;V_n+\sigma_d\,|I|\,R_n
\qquad(\sigma_d=+1\ \text{방전})
\label{eq:vapppol}
\end{equation}
로 적는다 --- 방전($\sigma_d=+1$)은 음극에서 산화(탈리튬화)가 일어나는 방향이라 측정 전위가 내부 전위보다 분극만큼
\emph{높게} 관측된다($V_\app>V_n$). 충전($\sigma_d=-1$)은 부호가 반대다. \textbf{(b) 연산 --- $V_n$ 으로 이항.}
식~\eqref{eq:vapppol} 을 $V_n$ 에 대해 풀면(분극 항을 좌변으로 옮겨)
\begin{equation}
V_n\;=\;V_\app-\sigma_d\,|I|\,R_n,
\label{eq:vnmid}
\end{equation}
\textbf{(c) 박스 --- 두 방향을 한 식으로.} $\sigma_d$ 가 두 방향을 모두 흡수하므로 식~\eqref{eq:vnmid} 가 곧 박스다:
\begin{equation}
\boxed{\;V_n\;=\;V_\app\;-\;\sigma_d\,|I|\,R_n\;.}
\label{eq:vn}
\end{equation}
부호 약속을 식으로 확인하면 ---
방전에서 $V_n=V_\app-|I|R_n<V_\app$(측정이 내부보다 높음), 충전에서 $V_n=V_\app+|I|R_n>V_\app$ --- 이며,
$|I|\to0$ 또는 $R_n=0$ 에서 $V_n=V_\app$ 로 두 전위가 일치한다.

\subsection{점별 평가 원칙 --- 자연 경계가 여는 여유}\label{sec:pointwise}
이후 모든 전이 계산은 내부 전위 $V_n$ 이 주어지는 낱낱의 점에서 직접 닫힌다 --- 전이 $j$ 의 기여
$[\text{peak shape}]_j(V_n)$ 은 그 점 하나만으로 완결되는 값이고, 공유된 보조 좌표를 거치지 않는다.
전이의 인과 기억(\S\ref{sec:tail} 의 식~\eqref{eq:lag})은 하한 $u\to-\infty$(방전) 또는 상한
$u\to+\infty$(충전)를 여는 적분이므로, 꼬리가 들어올 여유는 그 극한 자체가 이미 무한히 열어 둔다 ---
평가점 $V$ 근방의 과거를 몇 배의 $L_{V,j}$ 만큼만 들여다보면 지수 가중치 $e^{-|V-u|/L_{V,j}}$ 가 스스로
그 바깥을 무시할 만큼 작게 만드므로(적분의 유효 지지폭은 $L_{V,j}$ 규모), 전이마다 필요한 과거의 폭도
전이마다 다르게 자연히 정해진다. 온도가 배열 $T(V)$(비등온)이면 \emph{점별 양}($U_j\cdot w_j\cdot\xi_\eq$)은 각 $V_n$ 점의 $T(V_n)$ 값을
그대로 쓰고, \emph{전이당 상수}(분기 gap$\cdot$지연 길이 --- N3$\cdot$N7)는 진행 구간 전체의 단일 평균
$T_\rep=\overline{T}$ 한 값으로 평가한다; 스칼라면 둘 다 $T$ 를 공통으로 쓴다. 배경도 같은 점에서
먼저 깔리고($\dd Q/\dd V|_{V_n}\leftarrow C_\bg(V_n)$) 그 위에 모든 전이 기여가 점별로 누적되므로
($\mathrm N9$), 관측 범위 밖으로 계산을 밀어 두는 인위적 조작이 필요 없다.

\begin{keybox}
\textbf{두 전위.} $V_\app$(측정)$\cdot$$V_n$(내부, 식~\eqref{eq:vn})만 남는다. 다음 절부터 전이 루프의 모든
물리량 --- 평형 중심$\cdot$분기$\cdot$폭$\cdot$평형 진행률$\cdot$지연 꼬리 --- 는 이 $V_n$ 위에서 바로
평가되고, 그 값이 곧 출력이다.
\end{keybox}

\begin{bgbox}
\textbf{측정 원리 배경 --- 입력 세 양의 측정 대응(색인).} 본문이 전제하는 세 입력은 각각 확립된 측정의
산물이며, 대응만 색인한다(기법 상세는 범위 밖): (i) 준평형 $U_\mathrm{oc}(\bar x)$ $\leftrightarrow$
GITT\cite{weppner_huggins1977}(저율 정전류 곡선의 분극 보정 식~\eqref{eq:vn} 은 그 연속 근사판), (ii) 엔트로피
계수 $\partial U_\mathrm{oc}/\partial T$ $\leftrightarrow$ 전위차 엔트로피법\cite{reynier2003,swiderska2019}(부호$\cdot$모양의
상전이 유형 판독 지도는 \cite{baek_pilon2022}), (iii) 가역 발열 $-IT\,\partial U_\mathrm{oc}/\partial T$
$\leftrightarrow$ 등온 열량. 모델 식 자체는 측정 방식과 무관하게 성립하며, 계산 절차$\cdot$실측 대응 상세는
Part T 방법론 절(\S\ref{sec:method}).
\end{bgbox}
% 자산(v1.0.21 신규): [V21-Q5-02 측정 원리 배경 bgbox — D21-2 top3 ③(GITT 원전·엔트로피법·해석 지도 대응, 스코프 한정 가드)]

\begin{srcbox}
\textbf{다리 --- 측정 원리 박스의 두 원전.} \emph{(i) GITT 와 식~\eqref{eq:vn}.} 본문이 입력으로 전제하는 준평형 $U_\oc(\bar x)$ 는
Weppner--Huggins 의 GITT\cite{weppner_huggins1977} 가 재는 그 양이고, 식~\eqref{eq:vn} 의 분극 보정
$V_n=V_\app-\sigma_d|I|R_n$ 은 같은 $U_\oc$ 로 가는 연속 근사판이다:
\begin{center}\footnotesize
\begin{tabular}{@{}ll@{}}
\hline
본문 & Weppner--Huggins\cite{weppner_huggins1977}(GITT) 대응 \\
\hline
목표량 준평형 $U_\oc(\bar x)$ & 펄스 후 완화된 정지 전위(준평형 OCV) \\
분극 보정항 $\sigma_d|I|R_n$ (식~\eqref{eq:vn}) & 대기 완화로 제거되는 IR$\cdot$과전압 \\
저율 연속 곡선(연속 근사) & 펄스--완화 이산 점열(정통 측정) \\
\hline
\end{tabular}
\end{center}
GITT 는 각 점 완전 완화로 분극을 $0$ 으로 만들고, 식~\eqref{eq:vn} 은 저율 곡선에서 분극을 상수-저항 $R_n$
모형($\sigma_d|I|R_n$)으로 근사 제거하므로($R_n$ 의 조성$\cdot$전류 의존 무시) GITT 만큼 완전하지 않은 연속 근사다.
엔트로피계수 유형 분류(균질 고용체$\cdot$이온 질서화$\cdot$두-상 공존)는 Baek--Pilon 해석 지도\cite{baek_pilon2022}(tier B)를
판독틀로 빌리며, 실제 부호$\cdot$크기는 세 분포 합(config$+$vib$+$elec, Part T)과 1차 문헌\cite{reynier2003,allart2018}
으로 확정한다(tier B 1차 병기 규약).
\end{srcbox}

이제 본론 결과 사슬의 평형식들이 어디서 오는지를 통계역학 첫 원리에서 놓는 Part 0 로 넘어간다.

% 자산: [A-008] [A-009] [A-010] [A-011] [A-012] [A-013] [A-014] [A-070] [A-071] [A-072] [A-073] [A-074]
